\subsection{测试小节二}

本节演示定理环境与数学公式的引用。

\begin{definition}{测试定义}{def:test}
    这是一个测试用的定义。
\end{definition}

\begin{theorem}{测试定理}{thm:test}
    这是一个测试用的定理。
\end{theorem}

\begin{lemma}{测试引理}{lem:test}
    这是测试用的引理。
\end{lemma}

\begin{proposition}{测试命题}{prop:test}
    这是测试用的命题。
\end{proposition}

\begin{corollary}{测试推论}{cor:test}
    这是测试用的推论。
\end{corollary}

\begin{remark}{测试注记}{rem:test}
    这是测试用的注记。
\end{remark}

\begin{example}{测试例}{ex:test}
    这是测试用的例子。
\end{example}

\begin{proof}
    见 \autoref{thm:test} 的证明。
\end{proof}

定理类引用测试：\autoref{def:test}、\autoref{thm:test}、\autoref{lem:test}、\autoref{prop:test}、\autoref{cor:test}、\autoref{rem:test}、\autoref{ex:test}。

\medskip

下面测试公式环境。

\begin{equation}
    E = mc^2
    \label{eq:test-1}
\end{equation}

\begin{align}
    (a+b)^2 &= a^2 + 2ab + b^2 \\
    (a-b)^2 &= a^2 - 2ab + b^2
    \label{eq:test-2}
\end{align}

\begin{gather}
    x + y = z \\
    \int_0^1 f(x)\,\mathrm{d}x = 1
    \label{eq:test-3}
\end{gather}

公式引用测试：\eqref{eq:test-1}、\eqref{eq:test-2}、\eqref{eq:test-3}。
